\chapter{Cerf theory and parameterised handle constructions version}
\label{chap:cerf-theory-and-parameterised-handle-constructions-version}

We have defined the pseudo-isotopy group $\mathcal{P}(X)$ and want to study $\pi_0\mathcal{P}(X)$. Cerf discovered a great correspondence: 
$$\pi_0\mathcal{P}\cong \pi_1(\mathcal{F},\mathcal{E})$$
where $\mathcal{F}=\mathcal{F}(X):=\{f: X\times I\to I| f=\text{standard projection on a neighborhood of }\partial(X\times I)\}$ and $\mathcal{E}=\mathcal{E}(X)\subset \mathcal{F}(X)$ contains all such $f$ that have no critical points at all. The explicit correspondence is: Given $f\in \mathcal{P}(X)$, let $p$ be the standard projection from $X\times I$ to $I$. Then $p\circ f\in \mathcal{E}(X)$ since $f$ is a diffeomorphism of $X\times I$. It is not hard to show that $\mathcal{F}(X)$ is contractible, therefore we can choose a path $\gamma=\{\gamma_t:t\in I\}$ in $\mathcal{F}$ with $\gamma_0=p, \gamma_1=p\circ f$ both in $\mathcal{E}$. Then $\gamma$ is the corresponding element in $\pi_1(\mathcal{F},\mathcal{E})$. Conversely, given $f_t: X\times I\to I$ starting at $f_0=p$ and ending at $f_1=q\in \mathcal{E}$, consider the gradient vector field $V_q$ of $q$ on $X\times I$, integrating $V_q$ gives a diffeomorphism on $X\times I: (x,t)\to \phi_{V_q}(x,0)(t)$.

Then we reduce the question to the study of \emph{one parameter family of functions} on $X\times I$, namely, $f_t: X\times I\to I$ with $f_0,f_1$ no critical points at all. Like what we did in Morse theory, we can perturb the path fixing $f_0,f_1$ such that for all but finitely many $t\in I$, $f_t$ is a Morse function with ordered, distinct critical values. At the exceptional points, there appears either an $(i,i+1)$-birth/death point (two critical points of index $i$ and $i+1$ have the same critical value, but by the ordering condition, they must cancel in one direction) or an $(i,i)$-crossing point (two non-degenerate critical points of same index $i$ with crossing critical values). The Cerf diagram records the changes of critical values in $f_t: X\times I$. An example is shown in \Cref{fig:general-cerf-diagram}.

\begin{figure}[!ht]
    \fcapside{
        \caption{This is a Cerf diagram with 2 eyes, each eye begin from the birth to the end. All births and ends happen in $g^{-1}_t(1/2)$. Except the four birth/death points, every $g_t$ is Morse. For example, $g_{2/3}$ has two non-degenerate critical points of index 2 and 3. There are two crossing points in the diagram. Every arrow denotes a handle slide between critical points with the same index.}
        \label{fig:general-cerf-diagram}
    }{
        \includegraphics[width=0.4\textwidth]{general-cerf-diagram.png}
    }
\end{figure}

In this paper we are trying to find a Cerf diagram for $f_\beta$ with only one eye of index 2 and 3. Then let $g_t, t\in I$ be the corresponding one parameter family of functions on $X\times I$ with $g_0=p, g_1=p\circ f_\beta$. If we successfully find such a family $g_t$ with such a Cerf diagram, then the only things happen here are as follows: For some small $\epsilon$, $g_\epsilon$ births a degenerate critical point(standard local model is $g_\epsilon=x^3-y_1^2-y_2^2+y_3^2+y_4^2$ near the origin $p=(0,0,0,0,0
)$ with $g_\epsilon(p)=1/2$), and for $t> \epsilon$, the critical point splits into two non-degenerate critical points of indices 2 and 3(standard local model is $g_{\epsilon+s}=x^3-sx-y_1^2-y_2^2+y_3^2+y_4^2$). Right after the birth point at $g_\epsilon^{-1}(1/2)$, the level set is $g_{\epsilon+s}^{-1}(1/2)=X\# S^2\times S^2$, with the dual sphere of the unique 2-handle being $*\times S^2$ and the attaching sphere of the unique 3-handle being $S^2\times *$. Then the two spheres isotope in $X\# S^2\times S^2$ during the interval $t\in (\epsilon, 1-\epsilon)$. They may intersect each other at more than one point during the isotopy, but finally, right before $t=1-\epsilon$, they return to the dual position, i.e. intersect at exactly one point. Then general Morse theory tells that the (2,3)-handle pair can be cancelled, so a death occurs at $t=1-\epsilon$. But via an isotopy, we can assume one of the spheres is fixed throughout. In short, the $\{g_t,t\in I\}\in \pi_1(\mathcal{F},\mathcal{E})$ is completely governed by an element in $\pi_1(\Emb(S^2\times D^2, X\# S^2\times S^2), \Emb_0(S^2\times D^2, X\# S^2\times S^2))$ where $\Emb_0(S^2\times D^2, X\# S^2\times S^2)$ denotes all the embedded $S^2\times D^2$ such that $S^2\times 0$ transversely intersect $*\times S^2$ at a single point.

A similar story holds for a Cerf diagram with only one eye of index 1 and 2. To be explicit, $\{g_t,t\in I\}\in \pi_1(\mathcal{F},\mathcal{E})$ is completely governed by an element in $\pi_1(\Emb(S^1\times D^3, X\# S^1\times S^3), \Emb_0(S^1\times D^3, X\# S^1\times S^3))$ with similar notation for $\Emb_0(S^1\times D^3, X\# S^1\times S^3)$.

In particular, whenever we have a loop of framed embedded $S^i$ in $X\# S^i\times S^{n-i}$, which means the element lies not only in the relative $\pi_1$, but also in $\pi_1(\Emb(S^i\times D^{n-i}, X\# S^i\times S^{n-i}),*)$, tracing the flow lines it is not hard to show that the resulting diffeomorphism on $X$ is just the following composition:
$$\pi_1(\Emb(S^i\times D^{n-i}, X\# S^i\times S^{n-i}),*)\to \pi_0\Diff(X\setminus \nu S^{n-i-1},\partial)\xrightarrow{\bigcup \id_{\nu S^{n-i-1}}}\pi_0\Diff(X,\partial)$$
where the first map is the isotopy extension at $t=1$ and removing $\nu S^i=\nu (S^i\times *)$ in $X\# S^i\times S^{n-i}$, which is diffeomorphic to $X\setminus \nu S^{n-i-1}$, where the $S^{n-i-1}$ denotes the dual sphere of the $(i+1)$-handle.

Just like \emph{``Building a manifold $X$ from a Morse function is equivalent to finding a handle decomposition for $X$''}, we have an equivalent version for pseudo-isotopy: \emph{Finding a one-parameter family of functions on $X\times I$ is equivalent to finding a one-parameter family of handle decompositions for $X\times I$}.

Thus if we have a loop $\mathcal{H}=\{H_t, t\in I|H_0=H_1\}$ of handle decompositions of $X\times I$ which starts and ends at a standard cancelling position which can be made null, by the above statement, this corresponds to an element in $\pi_1(\mathcal{F},\mathcal{E})=\pi_0\mathcal{P}$, the resulting manifold is $Z_t$ which is diffeomorphic to $X\times I$. Since $H_0=H_1$, we have a natural way to identify $Z_0$ with $Z_1$ (since they are totally the same manifold). Then we build a cobordism from $S^1\times X$ to $Y=X\times I/((x,0)\sim (f_\mathcal{H}(x), 1))$ where $f_\mathcal{H}$ is the resulting diffeomorphism on $X$ with respect to $\mathcal{H}\in \pi_1(\mathcal{F},\mathcal{E})$. In handlebody language, we describe what $f_\mathcal{H}$ is: Each $H_t$ gives a series of surgeries $X=X_{0,t}\to X_{1,t}\to ...\to X_{n,t}$, $X_{n,t}$ is diffeomorphic to $X$ but not in a natural way. But again, since $H_0=H_1$, we have $X_{n,0}=X_{n,1}\cong X$. But as $t$ varies, since $H_t$ moves smoothly, every attaching region moves smoothly, thus we have a diffeomorphism $\phi_{m,t}: X_{m,0}\to X_{m,t}, \forall m,t$. In particular $\phi_{n,1}=f_\mathcal{H}:X=X_{n,0}\to X_{n,1}=X$ is the desired diffeomorphism on $X$.

In particular, when $\mathcal{H}$ is a loop of a single $(i,i+1)$-cancelling pair, that is, $Z_t=X\times I\cup h_{i,t}\cup h_{i+1,t}$. Since $h_{i,0}$ and $h_{i+1,0}$ form a standard cancelling pair, the attaching sphere of $h_{i,0}$ must be unknotted. Then let $X_{1,0}=(X\setminus  \nu S^{i-1}) \cup D^i\times S^{n-i}=X\# S^i\times S^{n-i}$ be the manifold after doing surgery on $h_i$. Up to an isotopy, we can assume $h_{i,t}=h_{i,0}$ throughout, so $X_{1,t}=X_{1,0}=X\# S^i\times S^{n-i}$. The only $t$-dependent data is the attaching region of $h_{i+1,t}$, which corresponds to $\pi_1(\Emb(S^i\times D^{n-i}, X_{1,0}),*)$, then one can directly see that the resulting $f_\mathcal{H}: X=X_{2,0}\to X_{2,1}=X$ is the composition map mentioned above.
